\section{Setting and reference predictors}
\paragraph{Families.} A system generates $x_t = a\,x_{t-1} + \sigma\varepsilon_t$ with $\varepsilon_t\sim\mathcal N(0,1)$ and
$x_0\sim\mathcal N(0,r^2)$ drawn independently of the system, \rlist. A context holds $x_0,\dots,x_n$ ($n$ transitions)
and the target is $x_{n+h}$. We use (a) a known noise scale ($\sigma=1$ in the maps; the first study also used a second
scale); (b) \emph{F3}: the same latent process observed as $y_t=x_t+v_t$, $v_t\sim\mathcal N(0,q)$, with noise ratios
$q\in\{0,\QOne,\QTwo,\QThree\}$, and a Kalman likelihood; (c) \emph{F1}: unknown noise scale; (d) \emph{F2}: a bivariate
VAR(1), $x_t=Ax_{t-1}+e_t$, $e_t\sim\mathcal N(0,\Sigma)$, with $x_0\sim\mathcal N(0,r^2 I)$. In every family the likelihood is
conditional on $x_0$, contexts are cut from the start of a trajectory by one function (one per observation model), and every
reference sees only the context.

\paragraph{References.} \emph{Population Bayes} (pop) is the Bayes predictive under the prior from which systems are drawn:
for (a)--(c), $a$ follows a normal distribution with mean \PriorMean and standard deviation \PriorSd truncated to
$|a|<\PriorRho$; in F1 the noise variance has an independent inverse-gamma prior with unit mean and unit variance; in F2 a
matrix-normal-inverse-Wishart prior chosen to mirror F1 (each noise variance has the same inverse-gamma marginal, and at unit
noise each coefficient has the same spread before truncation), truncated to spectral radius below \PriorRho.
\emph{Context-only Bayes} (stab in (a) and (b), ctx in F1 and F2) knows stability but nothing about the population. In (a)
and (b) it is a uniform prior on $a$ over $(-\StabBound,\StabBound)$, slightly wider than the population's support, and it is given
the true noise model (the noise scale, and the observation-noise variance in F3), as is the population reference: the two
differ only in their prior over $a$, so the gap isolates population knowledge about the dynamics. In F1 it adds a
log-uniform prior on $\sigma$ over $[\SigLo,\SigHi]$. In F2 it is a conjugate prior with zero-mean, very diffuse
coefficients (row covariance \FtwoCtxVzero times the identity) and an inverse-Wishart noise prior with \FtwoCtxNu degrees of
freedom and scale \FtwoCtxPsi times the identity, truncated to spectral radius below \StabBound; a conjugate prior cannot be
scale-free, so this comparator knows little, but not nothing, about the noise level. In F1 and F2 the gap therefore combines
knowledge of the dynamics and of the noise level (Section~\ref{sec:theory}). \emph{Library Bayes} (lib$_K$) is the posterior
over $K$ systems drawn from the population, with uniform weights; in F2 a library system has seven free parameters (four
coefficients and three noise-covariance entries). Matched priors are enforced by fingerprints: an optimality diagnostic
refuses a reference whose prior differs from the data's.

\paragraph{Computation.} Closed forms are used where they exist (conjugate updates, Student-$t$ predictives, library
mixtures). Integrals over $a$ use composite Gauss--Legendre quadrature with adaptive panels; in F2 the stability truncation
is handled by Monte Carlo over the posterior of $A$ with the noise covariance integrated analytically, and a control
variate equal to the closed-form untruncated predictive. Every approximation reports its own measured error
(Section~\ref{sec:val}).
